\paragraph{Articulation des coordonnées régionales et de l'incidence de clôture.}
\label{inc-ampl-articulacion}
La fonction de ce prolongement se comprend à partir d'une distinction
entre trois opérations : conserver un registre, évaluer une coordonnée et
extraire une relation d'incidence. Toutes trois agissent sur une information
produite par APP--TRIT--TPK. Leur articulation explicite la thèse
structurelle de l'article : une réalisation numérique possède une provenance
opératorielle, et certaines données de cette provenance admettent en outre une
réalisation combinatoire et réticulaire. Les équations permettent de suivre les deux
parcours avec leurs domaines et avec l'information qui subsiste dans chacun.

Le point de départ est l'état terminal muni de ses prolongements
compatibles. Il conserve les restes et quotients des deux feuillets
APP, le régime et l'orientation TRIT, ainsi que le transport TPK avec sa mémoire.
Les régions de clôture, de propagation et d'auto-échelle se déterminent dans ce
domaine avant que leurs lecteurs établissent les coordonnées
$\pi,e,\varphi$. Simultanément, le registre orienté des douze
fenêtres détermine le vecteur signé $U$. L'apparition ultérieure de $K$
provient de la transformation entière démontrée dans
\S\ref{subsec:k-hadamard} :
\begin{equation}
 U=\mathcal H_{12}K,\qquad
 K=\frac14\mathcal H_{12}U,\qquad
 \mathcal H_{12}^{2}=4I_{12}.
 \label{inc-ampl-registro}
\end{equation}
La congruence qui définit $\mathcal L_H$ garantit l'intégralité de
l'inverse. Ainsi, $K$ constitue une mise en coordonnées réversible
du registre signé dans le sous-réseau pertinent. Ses positions conservent
l'ordre de phase ; ses différences $D_3K,D_4K$ conservent les relations
transversales ; la somme $Q(K)$ détermine la composante uniforme. Cette
structure explique pourquoi le registre intervient à la fois dans la composition
numérique et dans la sélection d'une configuration de frontière.

L'évaluation périodique
\[
 K\longmapsto\kappa_{\mathrm{per}}
 =\frac{\sum_{m=1}^{12}K_m1000^{12-m}}{1000^{12}-1}
\]
maintient les douze blocs récupérables lorsque la base, la longueur et
l'origine de phase sont fixées. L'extraction d'un support réalise une opération
différente : elle conserve l'appartenance de chaque position à une classe définie
par un seuil et regroupe les registres qui possèdent le même support. La
distinction entre ces deux applications est décisive. La précision
scalaire de $\kappa_{\mathrm{per}}$ et l'incidence de $B_K$ sont deux
contenus mathématiques spécifiques, reliés par le registre
complet dont ils proviennent.

